% Self-contained Methods, Results, evidence table, and caption for Figure 4.
\subsection{Contrast-dependent hollowing and outer-shell dominance}
\label{sec:hollowing_contrast}

\paragraph{Methods.}
To compare hollowing across refractive-index contrast, we analyzed two 120-domain quarter-wave Bragg onions with air-like B domains and material A as the outermost domain: the $\lambda_0=470$~nm design with $n_A/n_B=1.55$ and the $\lambda_0=580$~nm design with $n_A/n_B=1.129$. Domains were replaced from the centre outward by air while preserving all original radial boundaries and the outer radius. Because every B domain already had $n_B=n_m=1$, only even hollowing depths were treated as physically distinct. Hollowing was therefore parameterized by the number of remaining high-index A domains, $N_{A,\mathrm{rem}}=(120-h)/2$. The principal retention metric was the wavelength-integrated backward-collected scattering relative to the intact sphere. Threshold shell counts were defined conservatively as the smallest $N_{A,\mathrm{rem}}$ reached continuously from the intact sphere without the retention curve first falling below the prescribed threshold. Peak retention and colour metrics were retained as supporting diagnostics, but the spectral integral was preferred because isolated resonant maxima were less robust to mode switching and numerical spikes.

\paragraph{Results.}
The number of outer high-index domains required to preserve the backward response depended strongly on refractive-index contrast. For the 470~nm design, 13 A domains were sufficient to retain at least 90\% of the intact backward spectral integral; two A domains retained at least 80\%, and one retained at least 50\%. The strongly hollowed five-domain geometry retained 84.9\% of the backward integral, although its largest backward peak was only 13.1\% of the intact maximum because the intact spectrum contained a much sharper resonance. The calculated chroma metric remained comparable to the intact structure, whereas the complete spectral line shape changed. For the lower-contrast 580~nm design, substantially more material was required: 31 A domains were needed for 90\% backward-integral retention, 19 for 80\%, and four for 50\%. With five A domains remaining, the backward peak and integral retentions were 55.0\% and 64.4\%, respectively, while the forward and total spectral integrals increased to 140.3\% and 136.7\% of the intact values. Thus, extensive hollowing can preserve the backward structural-colour response in the higher-contrast design, but the same degree of hollowing converts the low-contrast response toward predominantly forward scattering.

\begin{table}[t]
\centering
\caption{Evidence supporting the contrast-dependent hollowing statements. Thresholds use the backward spectral integral and the contiguous-from-intact definition described in the Methods.}
\label{tab:hollowing_thresholds}
\begin{tabular}{lcccc}
\hline
Design & $n_A/n_B$ & 90\% threshold & 80\% threshold & 50\% threshold\\
 & & \multicolumn{3}{c}{minimum remaining A domains}\\
\hline
470~nm & 1.55  & 13 & 2  & 1\\
580~nm & 1.129 & 31 & 19 & 4\\
\hline
\end{tabular}
\end{table}

\paragraph{Figure caption.}
\textbf{Contrast-dependent hollowing and outer-shell dominance.}
(a) Retention of the wavelength-integrated backward-collected scattering as high-index A domains are removed from the centre outward. Horizontal lines mark 90\%, 80\%, and 50\% retention. (b) Forward- and backward-integral changes relative to the intact spheres. Excessive hollowing of the low-contrast design increases forward scattering even as its backward response declines. (c,d) Selected backward-collected spectra for the 470 and 580~nm designs, including the intact sphere, the 90\%-retention shell count, a five-A-domain structure, and the all-air control. Numerical screening was restricted to nonfinite values and passivity violations in nominally lossless calculations; only isolated bounded gaps were interpolated in plotting copies, with every affected wavelength recorded in the accompanying audit table.
